% Insert after defining Q_{A,B} and its four templates, before the new cover.
% Uses the notation A,B,g and labels in truncated-square-section.tex.
\subsection{The Viterbo inequality for all partners}

The covering description also proves the Viterbo inequality for the entire
corner-truncated-square class, independently of which cover has minimum area.

\begin{theorem}\label{thm:truncated-global-viterbo}
Let $Q\subset\mathbb R^2$ be a parallelogram with one corner cut off by a
line meeting the interiors of its two adjacent sides. For every planar
convex body $K$ with nonempty interior,
\[
 c_{\mathrm{EHZ}}(K\times_L Q)^2<2|K|\,|Q|.
\]
In the normalization $Q=Q_{A,B}$ above, every capacity-one cover satisfies
\begin{equation}\label{eq:truncated-global-area}
 |K|>\frac{AB}{2AB-(A+B-1)^2}.
\end{equation}
Thus the strict-inequality assertion of \cite[Conjecture~6.1]{BMP2026}
holds, although its proposed optimizer will be disproved below.
\end{theorem}

\begin{proof}
An affine change in the $Q$-plane, paired with its inverse transpose in
the $K$-plane, reduces to $Q_{A,B}$ without changing the capacity or the
product of areas. Scaling $K$, we may assume its product capacity is one.
It therefore contains translated copies of the four templates
$T,-T,H,V$ from \eqref{eq:four-truncated-templates}.

First, this implies a lattice covering property. Put
\[
 \Lambda=A\mathbb Z e_1+B\mathbb Z e_2.
\]
Split every fundamental rectangle by its diagonal parallel to $(A,B)$.
The resulting triangulation has cells that are translates of $T$ or $-T$.
We use the triangulation-covering lemma
\cite[Lemma~4.3]{BMP2026}: if every cell of a planar triangulation fits by
translation into a convex body, then every translate of that body meets
the vertex set of the triangulation. It follows that
\begin{equation}\label{eq:truncated-lattice-cover}
 K+\Lambda=\mathbb R^2.
\end{equation}

Suppose first that $K$ contains a translate of the axis-parallel rectangle
$[0,A]\times[0,B]$. Let $w_x,w_y$ be the coordinate widths of $K$.
The four caps based on the rectangle's sides and reaching support-maximizing
points of $K$ have pairwise disjoint interiors
\cite[Lemma~7.1]{BMP2026}. Adding their areas to the rectangle's area gives
\[
 |K|\geq AB+\frac B2(w_x-A)+\frac A2(w_y-B)
       =\frac{Bw_x+Aw_y}{2}\geq\frac{A+B}{2},
\]
because $H,V$ fit into $K$ and hence $w_x,w_y\geq1$.
The cap lemma applies to an arbitrary convex body: only the four support
maximizers and the rectangle are used, so one may equivalently apply its
polygonal statement to their convex hull. This branch is strictly stronger
than the desired bound, since
\[
 \frac{A+B}{2}-\frac{AB}{2AB-g^2}
 =\frac{g\bigl(A(1-A)+B(1-B)\bigr)}{2(2AB-g^2)}>0.
\]

Now suppose no translate of that rectangle is contained in $K$.
For each $t\in\mathbb R^2$, form a finite graph whose vertices are
$(K-t)\cap\Lambda$, joining two vertices when their difference is
$Ae_1$ or $Be_2$. It is nonempty by \eqref{eq:truncated-lattice-cover}.
It is also a forest. Indeed, a cycle contains a simple rectilinear cycle;
its bounded interior contains a fundamental lattice cell. The closure of
that cell lies in the convex hull of the cycle vertices, hence in $K-t$.
This would put a translate of $[0,A]\times[0,B]$ inside $K$, contrary
to the present assumption. Thus, writing $N(t)$ for the number of vertices
and $E_x(t),E_y(t)$ for the two types of edges,
\[
 E_x(t)+E_y(t)\leq N(t)-1.
\]
Integrate this inequality over a fundamental rectangle of $\Lambda$.
Periodization of indicator functions gives
\[
 \int N(t)\,dt=|K|,\qquad
 \int E_x(t)\,dt=|K\cap(K+Ae_1)|,\qquad
 \int E_y(t)\,dt=|K\cap(K+Be_2)|.
\]
The fundamental rectangle has area $AB$, so
\begin{equation}\label{eq:forest-overlap}
 |K|-AB\geq |K\cap(K+Ae_1)|+|K\cap(K+Be_2)|.
\end{equation}

Finally, let $[u,u+e_1]\subset K$ be a placed copy of $H$.
Convexity gives the homothetic inclusion
\[
 (1-A)K+A(u+e_1)\subset K\cap(K+Ae_1),
\]
because subtracting $Ae_1$ from the left-hand side gives
$(1-A)K+Au\subset K$. Consequently
\[
 |K\cap(K+Ae_1)|\geq(1-A)^2|K|.
\]
The placed copy of $V$ gives the analogous bound
$|K\cap(K+Be_2)|\geq(1-B)^2|K|$. Inserting these inequalities
into \eqref{eq:forest-overlap} yields
\[
 \bigl(1-(1-A)^2-(1-B)^2\bigr)|K|\geq AB.
\]
The coefficient on the left equals $2AB-g^2>0$, which proves
the weak form of \eqref{eq:truncated-global-area}. Together with
$|Q_{A,B}|=(2AB-g^2)/(2AB)$, this proves the weak capacity-volume
inequality for every partner.

It remains to exclude equality. Suppose $|K|=AB/(2AB-g^2)$.
Replace $K$ by the convex hull $P$ of one placed copy of each of the four
templates. This polygon has at most ten vertices, is contained in $K$,
and still has product capacity at least one. The weak inequality already
proved gives $|P|\geq AB/(2AB-g^2)=|K|$, hence $P=K$.
Thus an equality body is a polygon. The rectangle case was strict, so
equality must hold in the two homothetic overlap inclusions. In particular,
with $[u,v]=[u,u+e_1]\subset K$,
\begin{equation}\label{eq:overlap-set-equality}
 K\cap(K+Ae_1)=(1-A)K+Av.
\end{equation}
The equality of sets follows from their inclusion and equal areas.

Let $n$ be an outward normal of any genuine facet of $K$.
The right side of \eqref{eq:overlap-set-equality} has a genuine facet
with that normal. An edge of an intersection of polygons comes from an
edge of one of the intersected polygons. Its supporting value is therefore
$h_K(n)+A\min\{0,n\cdot e_1\}$. Comparing supporting values in
\eqref{eq:overlap-set-equality} gives
\[
 (1-A)h_K(n)+A\,n\cdot v
 =h_K(n)+A\min\{0,n\cdot e_1\},
\]
so
\[
 h_K(n)=\max\{n\cdot u,n\cdot v\}.
\]
Every facet of $K$ consequently contains at least one endpoint of the unit
horizontal segment. Each endpoint can belong to at most two genuine
facets of a planar polygon. Hence $K$ has at most four facets.

If $K$ has four facets, their coverage by the two endpoints forces $u,v$
to be opposite vertices. Applying the same argument to the vertical
overlap equality shows that the endpoints of the unit vertical segment
are the other two opposite vertices. The two diagonals are perpendicular
and both have length one, so $|K|=1/2$. This is impossible because
$AB/(2AB-g^2)>1/2$.

The only remaining case is a triangular $K$. After translation, write
\[
 K=M\operatorname{conv}\{0,e_1,e_2\},\qquad \det M\ne0.
\]
The product symmetry $c(K\times_L Q)=c(Q\times_L K)$, followed by the
inverse-transpose linear change in the other factor, shows that $M^TQ$
is a capacity-one partner of the standard triangle. Its two normal
triangles triangulate the unit square grid. Applying the same
triangulation-covering lemma therefore gives
\[
 Q+M^{-T}\mathbb Z^2=\mathbb R^2.
\]
This lattice has covolume
\[
 |\det M|^{-1}=\frac{1}{2|K|}=|Q|.
\]
A lattice cover at this exact area has multiplicity one almost everywhere:
its multiplicity is at least one and its integral over a fundamental cell
is the cell's area. In particular distinct translates have disjoint
interiors, so these translates tile the plane.

Every edge of a convex polygon that tiles by lattice translations has
another edge parallel to it with the opposite outward normal. To see this,
choose a relative-interior point of the edge which is not a vertex of any
of the locally finitely many neighboring tiles. A tile covering points
just beyond that edge must meet it in an edge with the opposite normal;
interior overlap is impossible. But the truncation edge of $Q_{A,B}$ has
no edge with the opposite normal. Thus $Q_{A,B}$ cannot be such a tile.
This contradiction excludes equality and proves the strict inequality.
\end{proof}

The argument separates two possibilities: a full lattice cell already
contained in $K$ gives a direct cap-area estimate, while its absence makes
the lattice-point graphs acyclic and converts lattice covering into an
upper bound on overlap. The unit normal segments supply the opposite,
homothetic lower bounds on those overlaps. Equality would force a partner
with at most four sides; the four-sided case has too little area, and the
triangular case would force the truncated square to tile by translations.
